\clearpage\begingroup
\SourceLeaf{150}\NoSourceHead
\begin{center}
\textbf{\Large\textls[120]{BRIEFWECHSEL.}}\par\bigskip
[1.]\par\medskip
\textbf{\textsc{Schumacher} an \textsc{Gauss}.\quad Altona, 10. Junius 1810.}\par\medskip
\rule{.75\linewidth}{.4pt}\par
[Briefwechsel zwischen \textsc{Gauss} und \textsc{Schumacher} I, S. 43.]\par
\rule{.75\linewidth}{.4pt}
\end{center}
\noindent\hspace*{\parindent}\{Schon wieder, werthester Freund! werden Sie mit einem Briefe heimgesucht! und sogar mit einem Briefe, der eine Bitte enthält. \textsc{Repsold} hat seine Schüsseln schon aus dem groben gedreht, und was selten bey ihm kommt, grade Zeit, die Arbeit zu beginnen. Er ist zu blöde, selbst zu bitten, glaubt aber, Sie würden vielleicht, wenn Sie seine Noth wüssten, ihm zu Hülfe kommen. Ich würde ihm gerne nach der \textsc{Klügel}schen Theorie in \textsc{Gilberts} Annal. Jahrg. 1810 Stück 3[*)], sein 8füssiges Objectiv berechnen, indessen traue ich den Formeln K.s nicht, der schon früher nach dem vollkommensten, noch das \textls[120]{allervollkommenste} (\textit{sic}) Objectiv angab. Ich habe keine Zeit, die Formeln zu prüfen und ganz die Dioptrik zu treiben. Sie sehen also, in welcher Noth wir sind, wenn Sie uns nicht zu Hülfe kommen. Ich glaube Ihnen schon in einem vorigen Briefe gemeldet zu haben, dass das Objectiv 8 Fuss Brennweite haben soll. Bey dem Glase dazu ist
\begingroup\renewcommand{\thefootnote}{}\footnotetext{[*) \textit{Angabe eines möglichst vollkommenen achromatischen Doppel-Objectivs u.s.w.} von dem Professor \textsc{Klügel} in Halle, Annalen der Physik 34, Leipzig 1810, S. 265.]}\endgroup
\begin{center}
\begingroup\small\setlength{\tabcolsep}{3pt}
\begin{tabular}{cccc l r c r}
mittler.&Brech.&Verh.&des&Kronglases&$1{,}5157$,&Zerstreuung&$0{,}0051$\\
»&»&»&»&Flintglases&$1{,}6109$,&»&$0{,}0090$
\end{tabular}\endgroup
\end{center}
\noindent . . . . . . . . . . . . . . . . . . . . .\}
\HistoricalFoot{XI1.}{19}
\clearpage\endgroup
